\documentclass[12pt,a4paper]{article}

\usepackage[spanish,es-nodecimaldot]{babel}
\usepackage[utf8]{inputenc}
\usepackage[T1]{fontenc}
\usepackage{lmodern}
\usepackage{amsmath,amssymb}
\usepackage{graphicx}
\usepackage{booktabs}
\usepackage{hyperref}
\usepackage{enumitem}
\usepackage{geometry}
\usepackage{listings}
\usepackage{xcolor}

\geometry{margin=2.5cm}

\lstset{
  basicstyle=\footnotesize\ttfamily,
  breaklines=true,
  frame=single,
}

\begin{document}
\section{Investigación técnica (Defensa TFG): arquitectura y pipeline operativo}

\subsection{1) Alcance de este documento}

Este informe \textbf{no} repite:

\begin{itemize}
  \item \verb|docs/Personal Research/data-structure-and-canonical-schema-research-report.md|
  \item \verb|docs/Personal Research/qrdqn-research-report.md|
\end{itemize}

Se centra en los demás componentes clave del sistema para defensa ante tribunal.

\subsection{2) Arquitectura general del proyecto}

El sistema está dividido en dos fases:

\begin{enumerate}
  \item \textbf{Fase~1 (offline entrenamiento/validación)}
    \begin{itemize}
      \item Carga y limpieza de dataset (\verb|src/load_cicids2017.py|)
      \item Entorno Gymnasium sobre dataset (secuencial) para decisiones binarias PERMIT/BLOCK
            (\verb|src/rl_defender_env.py|)
      \item Entrenamiento (\verb|src/train_rl_defender.py|)
      \item Validación A/B/C y leave-one-exact-CSV-out (\verb|src/validate_checks.py|,
            \verb|src/validate_leave_one_csv_out.py|)
    \end{itemize}
  \item \textbf{Fase~2 (offline inferencia en tráfico real de laboratorio)}
    \begin{itemize}
      \item Entrada mantenida: \verb|scripts/predict_real_traffic_v2.py|
      \item Pipeline robusto con mapeo canónico, escalado, clipping y diagnósticos de shift.
    \end{itemize}
\end{enumerate}

\subsection{3) Procesado de datos (CICIDS2017) y anti-leakage}

\subsubsection{3.1 Configuración de carga}

La dataclass \texttt{CICIDSLoadConfig} controla rutas y preprocesado:

\begin{itemize}
  \item \texttt{chunksize=250\_000}
  \item \texttt{max\_rows}, \texttt{sample\_frac}
  \item \texttt{drop\_identifier\_cols=True}
  \item \texttt{scale=True}
  \item \texttt{use\_canonical=True}
  \item \texttt{test\_size=0.2}, \texttt{random\_state=42}
\end{itemize}

Referencia: \verb|src/load_cicids2017.py| (clase \texttt{CICIDSLoadConfig}, líneas $\sim$41--62).

\subsubsection{3.2 Limpieza y saneado}

Pipeline en \verb|_prepare_cicids_features(...)|:

\begin{enumerate}
  \item Detecta columna de etiqueta (\verb|_find_label_column|)
  \item Convierte etiqueta a binario (\texttt{BENIGN} $\to$ \texttt{0}, resto $\to$ \texttt{1})
  \item Elimina columnas con riesgo de leakage (\verb|_drop_identifier_like_columns|):
        Flow ID, Timestamp, IPs, puertos, etc.
  \item Fuerza numéricos (\verb|_coerce_numeric_features|)
  \item Sustituye \texttt{inf/-inf} por \texttt{NaN} y luego \texttt{NaN} $\to$ \texttt{0} en
        features (\verb|_clean_rows|)
  \item Elimina filas sin etiqueta
\end{enumerate}

Referencias: \verb|src/load_cicids2017.py| ($\sim$104--161, $\sim$221--272).

\subsubsection{3.3 Imputación y máscara de missingness}

La imputación explícita es \textbf{relleno con 0} para features de flujo (\texttt{fillna(0)}).
La semántica de máscara utilizada por el proyecto es:

\begin{itemize}
  \item \texttt{1 = presente / válido}
  \item \texttt{0 = ausente / imputado}
\end{itemize}

Referencias:

\begin{itemize}
  \item imputación: \verb|src/load_cicids2017.py| ($\sim$150--160)
  \item semántica: \verb|scripts/predict_real_traffic_v2.py|
        (\texttt{config["mask\_semantics"] = "1=present,0=missing"})
\end{itemize}

\begin{quote}
Nota para defensa: aquí conviene explicar que la imputación numérica y la máscara separan ``valor''
de ``confianza de disponibilidad''.
\end{quote}

\subsection{4) Entorno de RL usado en entrenamiento/validación}

Clase: \texttt{RLDatasetDefenderEnv} (\verb|src/rl_defender_env.py|).

\subsubsection{4.1 Interfaces Gymnasium}

\begin{itemize}
  \item \texttt{reset(...)} devuelve \texttt{(obs, info)}
  \item \texttt{step(action)} devuelve \texttt{(obs, reward, terminated, truncated, info)}
  \item \texttt{action\_space = Discrete(2)} (0=PERMIT, 1=BLOCK)
  \item \texttt{observation\_space = Box(shape=(n\_features,), dtype=float32)}
\end{itemize}

\subsubsection{4.2 Recompensa}

Config por defecto:

\begin{itemize}
  \item \texttt{tp=1.5}
  \item \texttt{fp=-1.5}
  \item \texttt{fn=-5.0}
  \item \texttt{omission=0.0}
\end{itemize}

La función \verb|_compute_reward(...)| penaliza fuertemente los falsos negativos (ataque permitido),
lo que prioriza seguridad.

\subsubsection{4.3 Dinámica del episodio}

\begin{itemize}
  \item Avance secuencial sobre muestras
  \item \texttt{shuffle=True} opcional en entrenamiento
  \item fin por \texttt{terminated} (fin dataset) o \texttt{truncated}
        (\texttt{max\_steps\_per\_episode})
\end{itemize}

Referencia: \verb|src/rl_defender_env.py| (constructor, \verb|_compute_reward|, \texttt{step}).

\subsection{5) Fase 2: inferencia robusta en tráfico real}

\subsubsection{5.1 Flujo operativo (main)}

\begin{enumerate}
  \item Lee CSV de flujos reales
  \item Separa metadatos (\texttt{src\_ip}, \texttt{dst\_ip}, \texttt{protocol}, puertos,
        \texttt{timestamp}, y columnas truth si existen)
  \item Armoniza unidades temporales (\verb|maybe_convert_time_units|): segundos $\to$
        microsegundos si detecta mediana $<1$
  \item Mapea al espacio canónico (\verb|map_to_canonical|)
  \item Aplica clipping percentilar opcional en features crudas (\texttt{--percentiles})
  \item Aplica escalado (\texttt{--scaler}) salvo \texttt{--no-scale}
  \item Aplica clipping z-score opcional (\texttt{--clip-z})
  \item Calcula diagnósticos de shift (\verb|compute_diagnostics|)
  \item Carga modelo (\verb|load_model|) y predice por lotes (\verb|batched_predict|)
  \item Exporta artefactos en \verb|runs/phase2/<RUN_ID>/|:
        \texttt{predictions.csv}, \texttt{config.json}, \texttt{metrics.json},
        \texttt{diagnostics.json} (si \texttt{--export-diagnostics})
\end{enumerate}

Referencia: \verb|scripts/predict_real_traffic_v2.py| ($\sim$141--475).

\subsubsection{5.2 Robustez ante outliers y drift}

Funciones de \verb|src/scaling_utils.py|:

\begin{itemize}
  \item \verb|apply_percentile_clipping(X, p_low, p_high)| (antes de escalar)
  \item \verb|apply_z_clipping(X_scaled, max_z)| (después de escalar)
\end{itemize}

Objetivo: evitar que outliers de tráfico real empujen al modelo fuera del régimen de entrenamiento.

\subsubsection{5.3 Métricas en inferencia con truth opcional}

\verb|compute_truth_metrics(...)| calcula, si hay columnas de verdad:

\begin{itemize}
  \item accuracy
  \item precision/recall/F1 de ataque
  \item TP/TN/FP/FN
\end{itemize}

Si no hay truth válida, la inferencia sigue y reporta métricas de distribución/predicción
(\texttt{block\_rate}, etc.).

\subsection{6) Mensaje clave para tribunal}

La arquitectura está pensada para:

\begin{itemize}
  \item entrenar y validar de forma controlada en Fase~1,
  \item ejecutar inferencia reproducible en Fase~2 con controles explícitos de drift (diagnósticos
        + clipping),
  \item manteniendo guardrails anti-leakage en la preparación de datos.
\end{itemize}

\end{document}
